\documentclass[12pt]{article}
\usepackage[utf8]{inputenc}
\usepackage[T1]{fontenc}
\usepackage[margin=1in]{geometry}
\usepackage{setspace}
\usepackage{xcolor}
\usepackage[colorlinks=true, urlcolor=blue!60!black, linkcolor=blue!60!black]{hyperref}
\setlength{\parindent}{0em}
\setlength{\parskip}{0.45em}
\renewcommand{\baselinestretch}{1.0}
\begin{document}
\begin{center}
{\Large \textit{Hammer v. Dagenhart}, 247 U.S. 251 (1918)\par}
\vspace{0.3em}
{\normalsize Case Brief --- One-Page Edition\par}
\end{center}

\textbf{Memory Jogger.} The child-labor case: manufacturing is not commerce, so Congress cannot ban shipment of harmless goods to regulate local labor; overruled by \textit{Darby} (1941) (\href{https://supreme.justia.com/cases/federal/us/247/251/}{247 U.S. 251}).

\textbf{Facts.} The Keating-Owen Act of 1916 barred interstate shipment of goods from mines and factories employing young children. Roland Dagenhart, whose two minor sons worked in a Charlotte cotton mill, sued U.S. Attorney W. C. Hammer to enjoin it.

\textbf{Procedural History.} The district court enjoined the Act; Hammer appealed directly.

\textbf{Votes.} Affirmed, 5--4. Day wrote for the Court; Holmes dissented, joined by McKenna, Brandeis, and Clarke.

\textbf{Holding.} Congress may not exclude intrinsically harmless goods from interstate commerce because of local labor conditions; doing so exceeds the commerce power and invades the police power reserved by the Tenth Amendment.

\textbf{Reasoning.} Earlier bans on lotteries, impure food, and liquor were valid because interstate transportation was necessary to the harm; here the cloth was harmless and the harm preceded shipment. Holmes answered that the Act regulated shipment itself, that ``regulation means the prohibition of something,'' and that indirect effects on state policy do not void a valid power.

\textbf{Examples.} \emph{Same-Side:} a minimum-wage furniture shipping ban falls. \emph{Opposite-Side:} bans on toxic toys or stolen cars stand. \emph{Fence-Sitter:} a ban on prison-made goods only into states that forbid them.

\textbf{Critique.} Textualists fault Day's insertion of ``expressly'' into the Tenth Amendment; progressives call the decision formalist and blind to the interstate race to the bottom.

\textbf{Key Quotation.} ``The making of goods and the mining of coal are not commerce, nor does the fact that these things are to be afterwards shipped or used in interstate commerce make their production a part thereof.'' --- Justice Day, Opinion of the Court, 247 U.S. at 272.

\end{document}
